%% Références citées dans la note. Les identifiants Sxx renvoient au registre
%% recherche/memo_stage/registre_sources_scientifiques.md, qui donne pour
%% chacune son fichier canonique, son empreinte et ce qu'elle établit
%% exactement — ainsi que les précautions de lecture.
\section{Références}
\label{sec:biblio}

\begingroup\emergencystretch=2em
\begin{description}[leftmargin=1.6em,style=nextline,itemsep=0.35em]

\item[Hendrick, Rinaldo et Manoli (2025)] Martin Hendrick, Andrea Rinaldo et
  Gabriele Manoli, « A stochastic theory of urban metabolism », \emph{PNAS},
  122(33), e2501224122. DOI \code{10.1073/pnas.2501224122}. \textsf{[S01]}\\
  \emph{Déclencheur de l'hypothèse directrice.} L'article établit un
  \emph{collapse} de distributions urbaines après remise à l'échelle. Il ne
  démontre ni auto-organisation critique, ni mécanisme de rebond : le passage
  à « une société économique pourrait être auto-critique » est une hypothèse
  formulée à partir de cette lecture.

\item[Bak, Chen, Scheinkman et Woodford (1992)] « Self-Organized Criticality
  and Fluctuations in Economics », Santa Fe Institute Working Paper
  1992-04-018. \textsf{[S02]}\\
  \emph{Antériorité à reconnaître.} Un modèle-jouet d'économie de production
  auto-critique existe depuis 1992. Ce travail ne revendique pas l'idée
  générale d'appliquer l'auto-organisation critique à l'économie.

\item[Drăgulescu et Yakovenko (2001)] « Exponential and power-law probability
  distributions of wealth and income in the United Kingdom and the United
  States », \emph{Physica A}, 299, 213--221. DOI
  \code{10.1016/S0378-4371(01)00298-9}. \textsf{[S03]}\\
  \emph{Cible distributionnelle initiale} : corps exponentiel « thermique »,
  queue de Pareto « superthermique ». Le corps exponentiel n'a jamais été
  obtenu dans ce modèle.

\item[Yakovenko et Rosser (2009)] « Colloquium: Statistical mechanics of money,
  wealth, and income », \emph{Reviews of Modern Physics}, 81, 1703--1725. DOI
  \code{10.1103/RevModPhys.81.1703}. \textsf{[S04]}\\
  Synthèse du programme précédent. Dans les données fiscales américaines
  qu'elle rassemble, la queue concerne environ 3~\% des déclarations en 1997 et
  son exposant varie dans le temps : \textbf{il n'existe pas de cible
  universelle} pour les revenus.

\item[Herbert, Giraud, Louis-Napoléon et Goupil (2022)] « Macroeconomic
  Dynamics in a Finite World: the Thermodynamic Potential Approach »,
  prépublication arXiv:2204.02038v2. \textsf{[S05]}\\
  \emph{Continuité théorique du postulat monnaie--exergie.} Ces auteurs fondent
  une macrodynamique de monde fini sur le couplage de sphères physique et
  économique \textbf{distinctes} : leur travail soutient la dépendance
  matérielle, \emph{pas} une identité monnaie--exergie.

\item[Dupont, Jeanmart et Possoz (2017)] « Vers un monde 100~\% renouvelable ? »,
  \emph{Regards économiques}, 135, p. 5--16. \textsf{[S06]}\\
  \emph{Illustration du couplage PIB--énergie} entre 1965 et 2015. La source
  mesure l'\textbf{énergie} et non l'exergie, et documente aussi un découplage
  relatif : elle illustre le postulat sans le valider.

\item[Ormerod et Mounfield (2001)] « Power law distribution of the duration and
  magnitude of recessions in capitalist economies: Breakdown of scaling »,
  \emph{Physica A}, 293(3--4), 573--582. DOI
  \code{10.1016/S0378-4371(01)00108-X}. \textsf{[S09]}\\
  \emph{Cible externe potentielle pour les récessions.} Elle impose de comparer
  les \textbf{distributions complètes} de durée et d'amplitude, et leurs
  ruptures éventuelles, pas seulement des moyennes. Cette comparaison n'a pas
  été faite : elle exige d'abord une correspondance temporelle.

\item[de Vries et Toda (2021)] « Capital and Labor Income Pareto Exponents
  across Time and Space », prépublication arXiv:2006.03441v3. \textsf{[S12]}\\
  \emph{Repère de comparaison le plus large disponible} : sur 475 observations
  pays--année couvrant 52 pays entre 1967 et 2018, les exposants du revenu du
  capital se situent principalement entre 1 et 3, contre 2 à 5 pour le revenu
  du travail. \textbf{Précaution obligatoire} : ces auteurs emploient
  l'exposant de fonction de survie, tandis que certains rapports de ce travail
  donnent un exposant de densité. La conversion doit être explicite avant toute
  comparaison.

\item[Sorrell et Dimitropoulos (2008)] Steve Sorrell et John Dimitropoulos,
  « The rebound effect: Microeconomic definitions, limitations and extensions »,
  \emph{Ecological Economics}, 65(3), 636--649.
  \href{https://doi.org/10.1016/j.ecolecon.2007.08.013}{DOI de l'article}.\\
  Distingue les définitions et hypothèses de mesure du rebond direct ;
  motive la séparation entre service utile, ressource et production du modèle.

\item[Wright (2005)] Ian Wright, « The duration of recessions follows an
  exponential not a power law », \emph{Physica A}, 345(3--4), 608--610.
  \href{https://doi.org/10.1016/j.physa.2004.07.035}{DOI : 10.1016/j.physa.2004.07.035}.
  Prépublication \href{https://arxiv.org/abs/cond-mat/0311585}{cond-mat/0311585}.
  \textsf{[S11]}\\
  Recompare les durées de récession d'Ormerod et Mounfield : l'exponentielle
  s'ajuste mieux aux données complètes de leur échantillon. Le traitement des
  épisodes d'un an change la comparaison. Ne pas confondre ce texte court avec
  l'autre article de Wright, \emph{The Social Architecture of Capitalism}.

\item[Reed (2003) ; Reed et Jorgensen (2004)] « The Pareto law of incomes---an
  explanation and an extension », \emph{Physica A}, 319, 469--486, DOI
  \code{10.1016/S0378-4371(02)01507-8} ; « The Double Pareto-Lognormal
  Distribution », \emph{Communications in Statistics}, 33(8), 1733--1753, DOI
  \code{10.1081/STA-120037438}.\\
  \emph{Loi de référence du processus de croissance-renouvellement} auquel M2
  se ramène. Son intérêt n'est pas d'être flexible à quatre paramètres : c'est
  la loi de l'état d'un mouvement brownien géométrique partant d'un état
  log-normal, observé après une durée exponentielle. C'est cette dérivation
  \emph{générative} qui la rend pertinente ici.

\item[Clauset, Shalizi et Newman (2009)] « Power-Law Distributions in
  Empirical Data », \emph{SIAM Review}, 51(4), 661--703.
  DOI \chemin{10.1137/070710111}.\\
  \emph{Méthode d'identification des lois de puissance} (maximum de
  vraisemblance et test de Kolmogorov--\allowbreak Smirnov) dont s'inspirent les scripts de
  test de queue de ce travail. C'est la source de la mise en garde
  « une régression log-log à bon $r^2$ ne suffit pas ».

\item[Wildauer et Heck (version septembre 2024)] Rafael Wildauer et Ines Heck,
  « Was Pareto right? Do the US wealth and income distributions follow power
  laws? », \emph{Greenwich Papers in Political Economy}, GPERC92
  (page de garde de série : 2023 ; version du texte : septembre 2024).
  \href{https://gala.gre.ac.uk/id/eprint/38597/13/38597\%20WILDAUER_Was_Pareto_right_Is_the_distribution_of\%20wealth_thick_tailed_(REVISED)_2023.pdf}{Archive institutionnelle de Greenwich}.\\
  Distingue Pareto strict et comportements asymptotiques ; les alternatives
  ne sont pas interchangeables. Le soutien à une queue asymptotiquement en
  puissance ne valide pas une loi de Pareto stricte sur tout le domaine.

\item[Shaikh, Papanikolaou et Wiener (2014)] Anwar Shaikh, Nikolaos
  Papanikolaou et Noe Wiener, « Race, gender and the econophysics of income
  distribution in the USA », \emph{Physica A}, 415, 54--60.
  \href{https://doi.org/10.1016/j.physa.2014.07.043}{DOI : 10.1016/j.physa.2014.07.043}.\\
  Étudie le programme exponentiel sur des sous-groupes de revenus du travail
  américains. Cette reprise ne démontre pas un consensus ni une loi universelle.

\item[Bouchaud (2024)] Jean-Philippe Bouchaud, « The Self-Organized Criticality
  Paradigm in Economics \& Finance », \href{https://arxiv.org/abs/2407.10284v2}{arXiv:2407.10284v2},
  6 septembre 2024. \textsf{[S07]}\\
  Distingue propagation critique et organisation endogène vers la criticité ;
  un rapport de branchement observé ne suffit pas au second énoncé.

\item[Dickman, Vespignani et Zapperi (1998)] Ronald Dickman, Alessandro
  Vespignani et Stefano Zapperi, « Self-organized criticality as an
  absorbing-state phase transition », \emph{Physical Review E}, 57(5),
  5095--5105. \href{https://doi.org/10.1103/PhysRevE.57.5095}{DOI : 10.1103/PhysRevE.57.5095}.
  \textsf{[S08]}\\
  Le substrat peut garder mémoire d'épisodes précédents ; la dynamique rapide
  d'une avalanche et l'évolution du système ne sont pas le même objet.

\item[Ribeiro et Rybski (2023)] Fabiano L. Ribeiro et Diego Rybski,
  « Mathematical models to explain the origin of urban scaling laws »,
  \emph{Physics Reports}, 1012, 1--39.
  \href{https://doi.org/10.1016/j.physrep.2023.02.002}{DOI : 10.1016/j.physrep.2023.02.002}.
  \textsf{[S10]}\\
  Plusieurs mécanismes produisent des lois d'échelle : la forme seule ne
  permet pas d'identifier le mécanisme de ce modèle.

\item[West, Brown et Enquist (1997)] Geoffrey B. West, James H. Brown et
  Brian J. Enquist, « A General Model for the Origin of Allometric Scaling
  Laws in Biology », \emph{Science}, 276(5309), 122--126.
  \href{https://doi.org/10.1126/science.276.5309.122}{DOI de l'article}.\\
  Modèle de transport en réseaux biologiques ; il éclaire la motivation
  allométrique sans imposer un exposant universel aux entités économiques.

\end{description}

Le registre complet, dans le dossier \chemin{recherche/memo_stage/}, fichier
\chemin{registre_sources_scientifiques.md}, donne pour
chaque source son fichier canonique, son empreinte, ce qu'elle établit
exactement et ce qu'elle \emph{n'établit pas}. Les entrées ci-dessus sont
identifiables sans consulter ce registre. Les dates de version des
prépublications sont conservées pour ne pas confondre versions et publications.
\endgroup
